\section{Storage as part of query execution}
\label{sec:pg-storage-role}

A biological sequence can be represented as a PostgreSQL variable-length
value, but this representation alone is insufficient for large reference
collections. A sequence may exceed the size of one PostgreSQL allocation. An
alignment may need only several bounded intervals from that sequence, while a
metadata query may not need sequence bytes at all. Storing every sequence as
one eagerly reconstructed value would therefore couple the cost of a query to
the total payload size even when the requested work is small.

BioSeqDB separates ordinary columns from externalized sequence payloads. The
table access method stores scalar values and sequence locators in a columnar
layout. Packed sequence fragments reside in relations selected by sequence
length and shard capacity. A locator in the parent row identifies the fragment
range and the immutable source version. Readers resolve that locator only when
an operator requests sequence data. Persistent minimizer indexes form a third
relation family, described in the next chapter. They are derived from the sequence
source and are not part of the primary payload layout.

This organization makes storage an execution property rather than a passive
serialization choice. Projection determines which column streams are read.
Scalar predicates can discard chunk groups before a sequence locator is
loaded. A lazy scanner bounds the fragment data retained by one execution.
PostgreSQL locks, snapshots, write-ahead logging (WAL), and publication rules
govern when newly written stripes and fragments become visible.

The columnar layer was adapted from Citus Columnar Storage.
BioSeqDB retains its page mapping \cite{citus_columnar_storage},
stripes, chunk groups, per-column streams, compression metadata, and min/max
pruning. It adds sequence-specific indirection and execution interfaces.
Table~\ref{tab:generic-columnar-comparison} summarizes the resulting
distinction.

\begin{table}[ht]
  \centering
  \small
  \caption{Generic column-store functions and BioSeqDB additions.}
  \label{tab:generic-columnar-comparison}
  \begin{tabularx}{\linewidth}{@{}p{0.27\linewidth}XX@{}}
    \toprule
    \textbf{Concern} & \textbf{Generic column store} &
    \textbf{BioSeqDB behavior}\\
    \midrule
    Physical organization & Stripes, chunk groups, compressed column streams &
    The same organization for scalar data and compact locators\\
    Large values & One serialized value, possibly managed by PostgreSQL TOAST &
    Fixed-size packed fragments addressed by a validated locator\\
    Filtering & Min/max pruning over ordered scalar values &
    Scalar pruning occurs before sequence access, while sequence values are
    excluded from min/max pruning\\
    Reading & Reconstruct projected rows from selected streams &
    Publish scalar or batch slots and acquire sequence fragments on demand\\
    Reuse & Backend-local decompression may avoid repeated decoding &
    Separate scopes for fragment residency, decompressed chunks, persistent
    indexes, and shared query results\\
    \bottomrule
  \end{tabularx}
\end{table}

\section{Columnar relation layout}
\label{sec:pg-storage-columnar}

\subsection{Logical byte space and reservations}

The main relation fork begins with two non-data blocks. Block~0 holds the
column-store metapage and block~1 is reserved. Logical stream data begin after
the PostgreSQL page header in block~2. Let
\[
  P = \mathrm{BLCKSZ} - \mathrm{SizeOfPageHeaderData}
\]
be the usable payload of one database page. Let \(x\geq 0\) be an offset
relative to the start of the stream-data area. It maps to
\[
  \mathrm{block}(x)=2+\left\lfloor\frac{x}{P}\right\rfloor,
  \qquad
  \mathrm{inpage}(x)=\mathrm{SizeOfPageHeaderData}+(x\bmod P).
\]
Higher layers can consequently reserve a continuous byte interval without
splitting every compressed stream at page boundaries.

The metapage records the format and storage identities together with the
first unused stripe number, row number, and byte offset. A writer reserves
non-overlapping ranges while holding the relation extension lock. It then
serializes its stripe into the reserved range. Metapage changes use
PostgreSQL's generic WAL support, so crash recovery follows the database page
and buffer model. Initialization and truncation require stronger relation
locks because they replace the shared allocation state.

A reservation is not equivalent to a visible stripe. Completion records the
actual row count, byte length, and stream metadata in the storage catalogs.
Readers use published metadata under their PostgreSQL snapshot. They do not
infer valid data from a byte range that a writer reserved but did not finish.
This distinction permits concurrent writers to allocate disjoint ranges while
preserving transactional visibility.

\subsection{Stripes, chunk groups, and streams}

A stripe is the large append unit. Each stripe is partitioned into chunk
groups, and every chunk group covers the same logical row interval for all
columns. The default configuration places up to 150,000 rows in a stripe and
10,000 rows in a chunk group. Both values are table options. Their purpose is
to balance metadata granularity, compression ratio, and the amount of data
handled in one scan step.

For each column in a chunk group, the writer produces two streams. An
existence bitmap represents NULL values at one bit per row. A value stream
contains aligned serialized values for present rows. The writer may compress
the value stream with Zstandard, LZ4, or PostgreSQL PGLZ, depending on the
build and table options. It retains the compressed form only when it is
smaller than the original. The chunk metadata records the algorithm that was
actually used, including the uncompressed case.

Figure~\ref{fig:column-store-layout} shows the hierarchy and its relationship
to a projected scan. The figure separates physical page mapping from the
logical units used by the reader.

\begin{figure}[ht]
  \centering
  \begin{tikzpicture}[x=0.9cm,y=0.9cm,font=\small,>=Latex]
    \node[draw,rounded corners,fill=blue!10,minimum width=3.0cm,
      minimum height=0.8cm] (meta) at (1.8,4.8) {block 0: metapage};
    \node[draw,rounded corners,fill=gray!12,minimum width=3.0cm,
      minimum height=0.8cm] (reserved) at (5.5,4.8) {block 1: reserved};
    \node[draw,rounded corners,fill=green!10,minimum width=4.0cm,
      minimum height=0.8cm] (bytes) at (10.0,4.8)
      {block 2 onward: stream bytes};

    \node[draw,rounded corners,fill=orange!12,minimum width=11.2cm,
      minimum height=0.9cm] (stripe) at (6.0,3.25)
      {stripe: reserved row range and byte range};
    \node[draw,fill=orange!22,minimum width=3.0cm,
      minimum height=0.75cm] (cg0) at (2.0,1.8) {chunk group 0};
    \node[draw,fill=orange!22,minimum width=3.0cm,
      minimum height=0.75cm] (cg1) at (6.0,1.8) {chunk group 1};
    \node[draw,fill=orange!22,minimum width=3.0cm,
      minimum height=0.75cm] (cgn) at (10.0,1.8) {later groups};
    \node[draw,rounded corners,fill=purple!10,minimum width=4.6cm,
      minimum height=0.85cm] (streams) at (4.0,0.35)
      {projected existence and values};
    \node[draw,rounded corners,fill=purple!10,minimum width=3.9cm,
      minimum height=0.85cm] (skip) at (9.4,0.35)
      {min/max pruning};

    \draw[->,thick] (bytes) -- (stripe);
    \draw[->,thick] (stripe) -- (cg0);
    \draw[->,thick] (stripe) -- (cg1);
    \draw[->,thick] (stripe) -- (cgn);
    \draw[->,thick] (cg0) -- (streams);
    \draw[->,thick] (cg1) -- (skip);
  \end{tikzpicture}
  \caption{Column-store hierarchy. Published metadata maps stripes and chunks
  to the logical byte space. Projection and pruning determine which streams
  are decoded.}
  \label{fig:column-store-layout}
\end{figure}

During a flush, offsets are calculated before any stream is copied. For one
column, the physical order is its existence buffers followed by its value
buffers. The completed metadata describes stripe and chunk-group row counts,
stream offsets, compressed and uncompressed lengths, deletion masks, and
compression identifiers. This metadata also supports columns added after an
older stripe was written. Such a column is reconstructed from its PostgreSQL
default or as NULL.

The metadata is divided by scope. A relation-options record contains the
stripe size, chunk-group size, compression method, and compression level. A
stripe record maps one storage identity and stripe number to its row interval
and logical byte interval. Chunk-group records contain row and deletion
counts. Per-column chunk records contain stream locations, value counts,
compression details, and optional extrema. Row-mask records represent deleted
logical rows without rewriting the immutable stream data. The sequence shard
registry is separate from these ordinary column-store records because its key
must also identify the parent column, length bucket, and fragment shard.
Keeping these scopes distinct allows a narrow metadata change to be recorded
without treating a fragment address as column-stream metadata.

The writer maintains min and max values for ordinary columns when PostgreSQL
provides an appropriate comparison procedure. These extrema support chunk
elimination. Sequence columns are excluded. A locator byte pattern has no
biological ordering, and the lexicographic order of packed symbols does not
provide a useful range predicate for alignment. A neighboring length, date,
or sample category remains eligible for pruning.

Several limits prevent a chunk from becoming an unbounded allocation. A
single aligned serialized value of 256~MiB or more is rejected. A stripe is
flushed before a value buffer approaches its approximately 1.024~GiB bound.
These limits apply to ordinary column chunks. Large biological sequences use
the fragment layout described next, so their logical length is not restricted
to one chunk buffer.

\section{Fragmented and packed sequence payloads}
\label{sec:pg-storage-sequence-indirection}

\subsection{Locator and fragment geometry}

An externalized sequence is represented in its parent row by a normal
\texttt{bioseq.seq} header followed by a 40-byte locator payload. The complete
stub occupies 56 bytes. The header marks the value as a locator while retaining
the alphabet identifier, packed storage width, and padding of the final
fragment. The current locator format is version~3. It records the format
version, length bucket, shard,
fragment count, logical sequence length, first fragment row, and source epoch.

The source epoch identifies an immutable content version. A row address alone
states where bytes were written, but it cannot show whether those bytes still
belong to the same logical sequence after a storage lifecycle change. Readers
therefore require a non-zero epoch and compare it with the registered source.
They also recompute the expected fragment count and final padding from the
logical length and packed representation. Invalid versions, reserved bytes,
geometry, or epochs are reported as data corruption rather than followed as
addresses.

The physical unit is alphabet-generic. For a packed storage unit of \(B\)
bytes that represents \(L\) letters, one fragment contains
\[
  U_f=1024\ \text{units}, \qquad
  S_f=1024B\ \text{bytes}, \qquad
  N_f=1024L\ \text{letters}.
\]
For DNA4 stored in 64-bit units, a full fragment occupies 8~KiB and represents
32,768 bases. Other alphabets change \(L\), and another packed unit changes
\(B\), without creating a DNA-only fragment path. Only the last fragment may
contain padding.

\subsection{Buckets, shards, and the canonical registry}

Known-size insertions select one of 35 length buckets from the bit width of the
final sequence length. The first bucket groups the smallest lengths, the
middle buckets distinguish bit widths 9 through 41, and the last bucket covers
larger values through the 64-bit range. A streaming insertion cannot know its
final size when writing starts, so it uses the largest bucket. Bucketing is an
operational choice rather than a biological classification. It prevents short
and very long values from competing for the same active shard.

A bucket can contain several shards. All fragments of one sequence occupy one
shard and receive contiguous row numbers. If \(r_0\) is the first row, the row
for fragment \(i\) is
\[
  r_i=r_0+i, \qquad 0\leq i< n_f.
\]
Known-size insertion estimates the required bytes and opens a new shard when
the active shard lacks sufficient configured capacity. Streaming insertion
checks the per-sequence cap while fragments arrive. A sequence is never split
across shards. If it exceeds the configured shard limit, insertion fails with
an error instead of creating a locator with ambiguous ownership.

The shard registry maps a parent storage identity, sequence column, bucket,
and shard to the corresponding fragment relation and attachment relation. Its
identity includes the database, tablespace, parent physical identity, column,
bucket, and shard. This scope prevents an address in one database or
tablespace from being accepted in another. The attachment relation can refer
to derived state, while canonical persistent minimizer generations remain
separate from the primary fragment rows.

Figure~\ref{fig:sequence-fragment-flow} follows one locator from the parent
row to a bounded algorithm request.

\begin{figure}[ht]
  \centering
  \begin{tikzpicture}[x=0.9cm,y=0.9cm,font=\small,>=Latex]
    \node[draw,rounded corners,fill=blue!10,minimum width=3.2cm,
      minimum height=1.0cm,align=center] (parent) at (1.7,3.4)
      {parent row\\56-byte locator};
    \node[draw,rounded corners,fill=gray!12,minimum width=3.5cm,
      minimum height=1.0cm,align=center] (registry) at (6.0,3.4)
      {shard registry\\bucket and shard};
    \node[draw,rounded corners,fill=green!10,minimum width=3.4cm,
      minimum height=2.4cm,align=center] (shard) at (10.5,3.4)
      {fragment shard\\row \(r_0\): fragment 0\\row \(r_0+1\): fragment 1\\later rows};
    \node[draw,rounded corners,fill=purple!10,minimum width=3.8cm,
      minimum height=1.0cm,align=center] (scanner) at (8.2,0.75)
      {lazy scanner\\bounded resident slots};
    \node[draw,rounded corners,fill=orange!12,minimum width=3.5cm,
      minimum height=1.0cm,align=center] (algorithm) at (3.3,0.75)
      {alignment stage\\requested interval};

    \draw[->,thick] (parent) -- (registry);
    \draw[->,thick] (registry) -- (shard);
    \draw[->,thick] (algorithm) -- (scanner);
    \draw[->,thick] (scanner) -- (shard);
    \draw[->,thick] (scanner) to[bend left=18] node[below] {packed view} (algorithm);
  \end{tikzpicture}
  \caption{Sequence indirection and bounded acquisition. The scanner resolves
  the validated locator and reads only fragments needed by the requested
  interval.}
  \label{fig:sequence-fragment-flow}
\end{figure}

\subsection{Insertion, concurrency, and cleanup}

Fragment allocation and writing use transaction-scoped advisory locks. One
key coordinates allocation for a parent column and bucket. A second key also
contains the shard and serializes writes to that internal resource. The keys
are compact coordination identifiers, not full catalog identities. A hash
collision can delay an unrelated writer, but it cannot make locators
interchangeable because subsequent catalog and epoch checks use the complete
identity.

Registry rows become authoritative only through the transaction that creates
or updates the storage resources. Byte deltas accumulated during writing are
applied when write states close. A reader holds a PostgreSQL snapshot and uses
normal relation ownership rules, so it cannot treat an uncommitted fragment
set as a complete source. Parent deletion and relation lifecycle handling are
also explicit. Fragment relations are registered storage resources rather
than permanent members of the extension installation. Parent-table drop
processing removes associated resources, and orphan cleanup covers mismatches
between registry rows and physical relations.

This separation matters because storage relations have a different lifetime
from SQL extension objects. It also avoids changing extension membership from
nested database execution contexts. The user-visible rule is simple. A
committed parent row either resolves to a validated fragment source or fails
with a storage error. It never silently reads a relation found only by a
reused physical address.

\section{Projection, pruning, and lazy reconstruction}
\label{sec:pg-storage-reading}

A scan begins with the requested columns and predicates. For each visible
stripe, the reader obtains its chunk metadata and determines which chunk
groups can satisfy eligible scalar predicates. It then loads only selected
streams for projected columns. The processing order is:

\begin{enumerate}
  \item read stripe and chunk metadata under the scan snapshot
  \item reject chunk groups contradicted by scalar min/max bounds
  \item read existence and value streams for the remaining projected chunks
  \item decompress when necessary and reconstruct aligned PostgreSQL values
  \item apply deletion masks and residual predicates
  \item resolve sequence locators only for rows and operators that need them
\end{enumerate}

Consider a relation containing a packed sequence, an identifier, collection
date, length, and sample attributes. A date predicate can reject a chunk group
without loading the sequence locator stream. A projection of identifiers and
lengths need not enter fragment storage at all. If an alignment operator needs
the sequence, it receives a lazy packed range rather than a concatenated
allocation.

The PostgreSQL-backed scanner retains a small set of open shard states for one
execution. Fragment data are copied into scanner-owned memory with 16 fixed
slots. With DNA4 and 64-bit units, the maximum full-fragment payload in those
slots is
\[
  16\times1024\times8 = 131{,}072\ \text{bytes}.
\]
A request contained in a resident fragment can return a direct packed view. A
request crossing a boundary is copied into bounded scratch space. Slot
replacement therefore decouples the resident sequence payload from the total
logical sequence length.

This bound applies only to fragment acquisition. Minimizer anchors, chaining
state, persistent index segments, and dynamic-programming matrices have their
own memory budgets. Fragmentation prevents an implicit whole-sequence copy,
but it does not make every alignment operation constant-memory.

The same column reader can publish scalar tuple slots or batch slots. A batch
contains up to one chunk group's rows, projected vectors, NULL masks, a keep
mask, and source row numbers. Supported comparisons and Boolean expressions
can update the keep mask before scalar row materialization. Supported metadata
aggregates may consume the batch interface. This use of the term vectorized
describes batched executor calls and does not imply that every transition is a
hardware SIMD instruction.

The planner uses the batch path only for a restricted single-relation plan
whose child can publish the required slot type. DISTINCT aggregates, aggregate
filters, alignment-owned projections, joins, and unsupported types retain the
ordinary scalar executor. Empty and all-NULL inputs for some integer and date
aggregates are a current limitation because their vector initial states do not
yet match PostgreSQL's NULL results. The optimization must therefore be
interpreted according to its stated eligibility and semantic limits.

Two local caches may reduce repeated work during these scans. Per-execution
fragment state keeps open shard relations and bounded fragment
slots. The decompression cache stores decoded compressed column chunks for one
PostgreSQL backend. Its key includes relation, stripe, chunk group, and column.
It has an explicit byte ceiling, skips chunks currently in use during
eviction, and returns borrowed buffers whose lifetime remains tied to the
cache entry. Neither mechanism coordinates work between sessions.

\section{Coordination of repeated query requests}
\label{sec:pg-storage-request-coordination}

Shared reuse of a completed SQL result is separate from fragment storage and
column decompression. The optional query-result cache can coalesce concurrent
executions with the same execution identity. Its identity includes the
analyzed query shape, parameter values, execution and security context, and a
dependency version vector. Normalized SQL text alone is insufficient because
two statements with the same text can run with different parameters or
visibility rules. The cache must be loaded when the PostgreSQL server starts
because its registry and memory areas are shared across backends.

The cache has three user-visible modes. In \texttt{disabled} mode, execution
does not create a cache claim. In \texttt{exact\_query} mode, the first
executor of an exact identity becomes the producer. Concurrent identical
executors wait, and later identical requests may replay the sealed ordered
multiset while its dependencies and lifetime remain valid. In
\texttt{reusable\_fragments} mode, exact coalescing remains available.
Broader reuse at an expensive operator boundary requires a provider that can
prove compatible coverage and retain the lineage needed for residual
predicates. Without that proof, behavior reduces to exact reuse. The current
BioSeqDB installation does not supply a production coverage provider, so
broader operator-level reuse remains unavailable.

The first exact claimant produces the result. A result becomes reusable only
after every cache row group and its manifest have been sealed and checked against the
schema and dependency versions. Producer failure or invalidation leaves no
partial result that another query can replay. A canceled waiter releases its
claim without canceling the producer, while a failed producer allows waiting
queries to run normally or establish a new generation.

Completed results use checked columnar cache row groups. Memory and file admission
have explicit ceilings. If neither tier can accept another row group, the
original query continues without publishing a reusable result. Expiration is
measured from successful publication, and cached results are cleared across a
server restart because they are derived state rather than durable table data.

\section{Resulting storage contract}
\label{sec:pg-storage-summary}

The storage design gives each layer a defined responsibility. The
Citus-derived column engine manages pages, stripes, chunk groups, scalar
projection, compression, deletion masks, and min/max pruning. Sequence
indirection adds alphabet-aware packed fragments, length buckets, shard
capacity, locator validation, and source epochs. Lazy scanners translate
bounded interval requests into fragment access. Shared query-result reuse has
its own identity, publication protocol, resource admission, and failure
behavior.

These boundaries support database execution without presenting a sequence as
an ordinary column-store value. Metadata filters can reduce the candidate set
before any sequence read. Alignment stages can request only the intervals they
need. Concurrent writers reserve independent physical ranges, while readers
observe only committed metadata and validated fragment sources. Repeated
queries may share a sealed result, but an incomplete or stale generation is
never exposed as complete. The result is a PostgreSQL storage model in which
logical sequence length, scan projection, and retained memory are related by
explicit contracts rather than by eager materialization.
